\subsection{Information Bounds and Asymptotic Efficiency}

Information inequalities describe the best possible local accuracy under regular models.  Efficient procedures are those whose asymptotic error reaches these bounds.

\begin{lemma} \label{lem:fisher-information-second-derivative} \hlt{Two Forms of Fisher Information}\\
For a regular one-parameter dominated model, let $\ell_\theta(x)\defeq\log p_\theta(x)$ and $S_\theta(X)\defeq\partial_\theta\ell_\theta(X)$.\\
Then $I(\theta)\defeq\mathbb{E}_\theta S_\theta(X)^2=-\mathbb{E}_\theta\partial_\theta^2\ell_\theta(X)$.
\begin{proof}
Differentiating $\int p_\theta\,d\mu=1$ gives $\mathbb{E}_\theta S_\theta=0$.\\
Differentiating, $0=\partial_\theta\mathbb{E}_\theta S_\theta=\mathbb{E}_\theta[\partial_\theta S_\theta+S_\theta^2]=\mathbb{E}_\theta[\partial_\theta^2\ell_\theta+S_\theta^2]$, since $\partial_\theta p_\theta=S_\theta p_\theta$. Thus $I(\theta)=\mathbb{E}_\theta S_\theta^2=-\mathbb{E}_\theta\partial_\theta^2\ell_\theta$.\\
\hlt{Structural Commentary}: Fisher information is both score variance and expected log-likelihood curvature.\\
\hlt{Key Insight}: Normalisation forces the mean score to vanish, linking slope variance to curvature.
\end{proof}
\end{lemma}

\begin{theorem} \label{thm:cramer-rao-bound} \hlt{Cramer--Rao Bound}\\
Let $X_1,\ldots,X_n\sim P_\theta$ be iid from a regular one-parameter dominated model with one-observation information $I(\theta)>0$. If $\widehat{\psi}\in L^2(P_\theta)$ is unbiased for differentiable $\psi(\theta)$, then $\operatorname{Var}_\theta(\widehat{\psi})\geq \frac{\psi'(\theta)^2}{nI(\theta)}$.
\begin{proof}
Let $S_n(\theta)\defeq\partial_\theta\log L_n(\theta)=\sum\limits_{i=1}^{n}S_\theta(X_i)$. Regularity gives $\mathbb{E}_\theta S_n=0$ and $\operatorname{Var}_\theta(S_n)=\mathbb{E}_\theta S_n^2=nI(\theta)$.\\
Differentiating $\mathbb{E}_\theta\widehat{\psi}=\psi(\theta)$ gives $\psi'(\theta)=\mathbb{E}_\theta[\widehat{\psi}S_n]=\operatorname{Cov}_\theta(\widehat{\psi},S_n)$.\\
Thus $\psi'(\theta)^2\leq\operatorname{Var}_\theta(\widehat{\psi})\operatorname{Var}_\theta(S_n)=\operatorname{Var}_\theta(\widehat{\psi})nI(\theta)$ by Cauchy--Schwarz.\\
\hlt{Structural Commentary}: The score is the direction in which unbiased estimators must covary to track $\psi(\theta)$.\\
\hlt{Key Insight}: Information lower-bounds variance through Cauchy--Schwarz.
\end{proof}
\end{theorem}

\begin{example} \label{ex:poisson-cramer-rao-attainment} \hlt{Poisson Mean Attains the Bound}\\
For $X\sim\operatorname{Poisson}(\lambda)$, $\ell_\lambda(X)=X\log\lambda-\lambda-\log(X!)$ and $\partial_\lambda^2\ell_\lambda(X)=-\frac{X}{\lambda^2}$, so $I(\lambda)=-\mathbb{E}_\lambda\partial_\lambda^2\ell_\lambda(X)=\frac{1}{\lambda}$.\\
For iid $X_1,\ldots,X_n$, any unbiased estimator $T$ of $\lambda$ satisfies $\operatorname{Var}_\lambda(T)\geq\frac{\lambda}{n}$.\\
$\overline{X}$ attains equality since $\mathbb{E}\overline{X}=\lambda$ and $\operatorname{Var}(\overline{X})=\frac{\lambda}{n}$.
\end{example}

\begin{theorem} \label{thm:mle-asymptotic-normality} \hlt{Regular MLE Asymptotic Normality}\\
Let $X_1,\ldots,X_n\sim P_{\theta_0}$ iid in a regular one-parameter model, with $\theta_0$ interior, $I(\theta_0)>0$, and a consistent interior MLE $\widehat{\theta}_n$. Then $\sqrt{n}(\widehat{\theta}_n-\theta_0)\Rightarrow N(0,I(\theta_0)^{-1})$.
\begin{proof}
Let $S_n(\theta)\defeq\partial_\theta\ell_n(\theta)$ and $J_n(\theta)\defeq-n^{-1}\partial_\theta S_n(\theta)$.\\
Since $S_n(\widehat{\theta}_n)=0$, Taylor expansion gives $0=S_n(\theta_0)-nJ_n(\widetilde{\theta}_n)(\widehat{\theta}_n-\theta_0)$ for some $\widetilde{\theta}_n$ between $\widehat{\theta}_n$ and $\theta_0$.\\
Hence $\sqrt{n}(\widehat{\theta}_n-\theta_0)=J_n(\widetilde{\theta}_n)^{-1}n^{-\frac{1}{2}}S_n(\theta_0)$.\\
Regularity gives $J_n(\widetilde{\theta}_n)\xrightarrow{P}I(\theta_0)$ and $n^{-\frac{1}{2}}S_n(\theta_0)\Rightarrow N(0,I(\theta_0))$, so Slutsky gives the result.\\
\hlt{Structural Commentary}: Local likelihood curvature converts score noise into estimator error.\\
\hlt{Key Insight}: The MLE is asymptotically score divided by information.
\end{proof}
\end{theorem}

\begin{definition} \label{def:asymptotic-efficiency} \hlt{Asymptotic Efficiency}\\
For a scalar regular parameter with information $I(\theta_0)>0$, an estimator sequence $\widehat{\theta}_n$ is \hlt{asymptotically efficient} at $\theta_0$ if $\sqrt{n}(\widehat{\theta}_n-\theta_0)\Rightarrow N(0,I(\theta_0)^{-1})$, i.e. it attains the inverse Fisher-information variance.
\end{definition}

\begin{remark} \label{rem:asymptotic-efficiency-versus-finite-sample-efficiency} \hlt{Efficiency and Approximation}\\
Finite-sample efficiency means equality in the Cramer--Rao bound.\\
Asymptotic efficiency means $\sqrt{n}(\widehat{\theta}_n-\theta_0)\Rightarrow N(0,I(\theta_0)^{-1})$; finite-sample bias or boundary effects still matter.
\end{remark}
